\section{The alpha-scaled ladder law}
\label{sec:ladder}

Consider a descending threshold ladder.  A rung has live threshold $\tau$,
the preceding threshold is $b\tau$, and an entering coordinate is assumed to
satisfy
\begin{equation}
 \tau<\frac{|r_u|}{\sqrt{d_u}}\leq b\tau.
 \label{eq:band-assumption}
\end{equation}
This in-band statement deliberately ignores fresh neighbor arrivals; that
limitation is audited in \cref{sec:spider}.

\begin{proposition}[Exact admissible bases]
\label{prop:admissible-bases}
For $\omega=\omega_\star$, the self-reflection from a push satisfying
\cref{eq:band-assumption} is guaranteed to lie below the current threshold if
and only if
\begin{equation}
 b\leq b_{\rm self}(\alpha):=\lambda_\alpha^{-2}.
 \label{eq:self-base}
\end{equation}
It is guaranteed to lie below the next threshold $\tau/b$ if and only if
\begin{equation}
 b\leq b_{\rm next}(\alpha):=\lambda_\alpha^{-1}.
 \label{eq:next-base}
\end{equation}
\end{proposition}

\begin{proof}
By \cref{lem:optimal-parameterization}, the worst in-band self-reflection has
scaled magnitude $\lambda_\alpha^2b\tau$.  Requiring this to be at most
$\tau$ gives $b\leq\lambda_\alpha^{-2}$; requiring it to be at most
$\tau/b$ gives $b^2\leq\lambda_\alpha^{-2}$.
\end{proof}

\begin{corollary}[Accelerated rung count]
\label{cor:rung-count}
Fix $s\in\{1,2\}$ and use $b_\alpha=\lambda_\alpha^{-s}$.  To cross a
threshold dynamic range $R>1$, the required number of rungs is
\begin{equation}
 K_s(R,\alpha)
 =\left\lceil
   \frac{\log R}
        {2s\operatorname{arctanh}(\sqrt\alpha)}
  \right\rceil
 =\Theta\!\left(\frac{\log R}{\sqrt\alpha}\right)
 \quad(\alpha\downarrow0).
 \label{eq:rung-count}
\end{equation}
\end{corollary}

\begin{proof}
The identity
$-\log\lambda_\alpha=2\operatorname{arctanh}(\sqrt\alpha)$ gives the exact
display.  For $0<t\leq1/2$, $t\leq\operatorname{arctanh}(t)
\leq t/(1-t^2)$, which proves the asymptotic order.
\end{proof}

\begin{proposition}[Fixed-base obstruction]
\label{prop:fixed-base}
For every fixed $b>1$, optimal SOR violates current-rung no-reactivation for
all
\begin{equation}
 \alpha<\left(\frac{\sqrt b-1}{\sqrt b+1}\right)^2.
 \label{eq:fixed-base-threshold}
\end{equation}
In particular, no fixed $b>1$ supplies an asymptotic no-self-reactivation
proof as $\alpha\downarrow0$.
\end{proposition}

\begin{proof}
Condition \cref{eq:self-base} is equivalent to
$\lambda_\alpha\leq b^{-1/2}$.  Substituting
$\lambda_\alpha=(1-\sqrt\alpha)/(1+\sqrt\alpha)$ and solving for
$\sqrt\alpha$ gives the complement of \cref{eq:fixed-base-threshold}.
\end{proof}

For calibration, $b_{\rm self}(0.04)=2.25$ and
$b_{\rm self}(0.025)\approx1.892$.  Thus the measured base-two prediction is
well aligned with the tested parameter range, but it is already slightly
outside the worst-case self-reactivation window at $\alpha=0.025$ and becomes
increasingly unsafe for smaller $\alpha$.
